\documentclass{beamer}
\usepackage[utf8]{inputenc}

\usetheme{Madrid}
\usecolortheme{default}
\usepackage{natbib}
\usepackage{makecell}
\usepackage{booktabs}
\usepackage{graphicx}
\usepackage{adjustbox}
\usepackage[dvipsnames]{xcolor}
\bibliographystyle{apalike}
\renewcommand{\cellalign}{cl}
\renewcommand{\theadalign}{cl}

% Where the auto-generated tables and figures live (relative to main.tex)
\graphicspath{{figures/}}

%------------------------------------------------------------
\title[]
{Expected Returns in \\International Government Bond Markets}

\subtitle{Master Thesis -- Mid-term Progress Presentation}

\author[Jan Heissenberger]
{Jan Heissenberger\\
Supervised by Giorgia Simion, Ph.D.}

\date[\today]
{\today}
%------------------------------------------------------------

\AtBeginSection[]
{
  \begin{frame}
    \frametitle{Table of Contents}
    \tableofcontents[currentsection]
  \end{frame}
}

\begin{document}

\frame{\titlepage}

\begin{frame}
\frametitle{Table of Contents}
\tableofcontents
\end{frame}

%==========================================================
\section{Recap}
%==========================================================

\begin{frame}{Recap -- the question}
\textbf{Main question.} Does the \cite{cieslak_expected_2015} cycle factor (CF) predict bond risk premia outside the US, and is its predictive content subsumed by a global cycle factor (GCF)?

\vspace{0.3cm}
\textbf{Subquestions}
\begin{enumerate}
    \item What is the role of the local CF relative to a GDP-weighted GCF? \\
    \color{gray}\small Integrated bond markets imply a shared risk factor across countries \citep{dahlquist_international_2013}.
    \color{black}
    \item Does the GCF subsume the local CF? \\
    \color{gray}\small Compensation for global risk should dominate local risk under integration.
    \color{black}
    \item Does an FX-adjusted global factor (FXGCF) improve return predictability for the USD investor?
\end{enumerate}
\end{frame}

\begin{frame}{Recap -- framework in one slide}
\textbf{Yield decomposition} \hfill (eq.\ 1--2)
\[
y_{i,t}^{(n)} = \alpha_{i,n} + \beta_{i,n}\,\pi^{e}_{i,t} + c_{i,t}^{(n)},\qquad
\pi^{e}_{i,t}=\tfrac{1-v}{1-v^{M}}\sum_{j=0}^{M-1} v^{j}\pi_{i,t-j}
\]

\textbf{Local cycle factor} \hfill (eq.\ 4--6)
\[
rx_{i,t+12} = \gamma_0 + \gamma_{i,1} c_{i,t}^{(1)} + \gamma_{i,2}\bar c_{i,t} + \varepsilon_{i,t+12},
\quad CF_{i,t}=\hat\gamma_{i,1} c_{i,t}^{(1)}+\hat\gamma_{i,2}\bar c_{i,t}
\]

\textbf{Global \& FX-global factors} \hfill (eq.\ 7, 15--17)
\[
GCF_t = \sum_{i\in G10} w_{i,t}\,CF_{i,t},\qquad
FXGCF_t = \hat\delta_0+\hat\delta_1\,GCF_t
\]
\centering
\small from a regression of GDP-weighted USD-investor returns on $GCF_t$.
\end{frame}

%==========================================================
\section{Milestones}
%==========================================================

\begin{frame}{Where we stand}
\begin{block}{Done}
\begin{itemize}
    \item \textbf{Data}: G10 zero yields, core CPI, FX vs USD, nominal GDP -- monthly panel, 1990s--2025.
    \item \textbf{Trend inflation}: DMA filter ($v=0.987$, $M=120$); cycle component as residual of $y\sim\pi^e$.
    \item \textbf{Returns}: maturity-averaged duration-standardized $rx$ (local) and $rx^{USD}$ (US investor).
    \item \textbf{Factors}: country-level CF, GDP-weighted GCF, FX-adjusted FXGCF, Cochrane--Piazzesi (2005) baseline -- with fully-OOS variants ($CF^{oos}, GCF^{oos}, FXGCF^{oos}$).
    \item \textbf{Inference}: per-country HAC (Newey--West, $L=18$) and panel Driscoll--Kraay SEs; nested Wald, Clark--West, Diebold--Mariano (in-sample \emph{and} on OOS factors).
    \item \textbf{Replication}: US Cieslak--Povala vs.\ Cochrane--Piazzesi -- in-sample and OOS $R^2$.
\end{itemize}
\end{block}
\begin{block}{In progress}
\begin{itemize}
    \item Robustness checks (sub-samples, alternative trend windows).
    \item Drafting Sections 4 (results) and 5 (discussion) of the thesis.
\end{itemize}
\end{block}
\end{frame}

\begin{frame}{Sample coverage}
\centering
\adjustbox{max width=\textwidth}{\input{tables/sample_coverage.tex}}

\vspace{0.3cm}
\small \textit{Estimation panel after dropping rows with any missing CF, GCF or FXGCF observation.}
\end{frame}

\begin{frame}{Yield panel and trend inflation}
\centering
\includegraphics[width=0.49\textwidth]{yield_ts.pdf}\hfill
\includegraphics[width=0.49\textwidth]{inflation_trend.pdf}\\
\small Left: zero yields by country/maturity. \quad
Right: YoY core CPI vs.\ DMA trend $\pi^e$.
\end{frame}

\begin{frame}{FX spot rates and nominal GDP}
\centering
\includegraphics[width=0.49\textwidth]{fx_ts.pdf}\hfill
\includegraphics[width=0.49\textwidth]{gdp_levels.pdf}\\
\small Left: monthly FX spot rates, USD per unit of foreign currency (log scale). \quad
Right: annual nominal GDP in USD, log scale.
\end{frame}

%==========================================================
\section{Results}
%==========================================================

\begin{frame}{Cycles and cycle factors}
\centering
\includegraphics[width=0.95\textwidth]{cycles_by_country.pdf}\\[-2pt]
\small Cycle component $c_{i,t}^{(n)}$ from $y\sim\pi^e$; cross-country co-movement is visible at long maturities.
\end{frame}

\begin{frame}{Local CF and Global CF}
\centering
\includegraphics[width=0.49\textwidth]{local_cf.pdf}\hfill
\includegraphics[width=0.49\textwidth]{gcf.pdf}\\
\small Left: $CF_{i,t}$ per country (eq.\ 6). \quad Right: GDP-weighted $GCF_t$ (eq.\ 7--8).
\end{frame}

\begin{frame}{Eq.\ (18) -- Local CF predicts local returns}
\centering
\scriptsize\textbf{In-sample (CF)}\\
\adjustbox{max width=\textwidth}{\input{tables/eq18_local_cf.tex}}

\vspace{0.15cm}
\scriptsize\textbf{Out-of-sample (CF$^{oos}$)}\\
\adjustbox{max width=\textwidth}{\input{tables/eq18_local_cf_oos.tex}}

\vspace{0.1cm}
\tiny HAC SEs (Newey--West, $L=\max(18,\,1.3\sqrt{T})$).\quad
$^{*}/^{**}/^{***}$ at the 10\%/5\%/1\% level.
\end{frame}

\begin{frame}{Eq.\ (20) -- Global CF predicts local returns}
\centering
\scriptsize\textbf{In-sample (GCF)}\\
\adjustbox{max width=\textwidth}{\input{tables/eq20_global_cf.tex}}

\vspace{0.15cm}
\scriptsize\textbf{Out-of-sample (GCF$^{oos}$)}\\
\adjustbox{max width=\textwidth}{\input{tables/eq20_global_cf_oos.tex}}
\end{frame}

\begin{frame}{Eq.\ (19) -- Does GCF subsume local CF?}
\centering
\scriptsize\textbf{In-sample (CF + GCF)}\\
\adjustbox{max width=\textwidth}{\input{tables/eq19_subsumption.tex}}

\vspace{0.15cm}
\scriptsize\textbf{Out-of-sample (CF$^{oos}$ + GCF$^{oos}$)}\\
\adjustbox{max width=\textwidth}{\input{tables/eq19_subsumption_oos.tex}}

\vspace{0.1cm}
\tiny Wald $p$: $H_0:\beta_{CF}=0$ in $rx_{i,t+12}=\alpha_i+\beta_i\,CF_{i,t}+\gamma_i\,GCF_t+\varepsilon$.
\end{frame}

\begin{frame}{Eq.\ (19) -- Coefficients with HAC bands}
\centering
\includegraphics[width=0.95\textwidth]{coef_eq19.pdf}\\
\small Per-country point estimates of $\beta_{CF}$ and $\gamma_{GCF}$ (eq.\ 19) with $\pm 1.96\,SE$.
\end{frame}

\begin{frame}{Eq.\ (22) -- USD investor on GCF}
\centering
\scriptsize\textbf{In-sample (GCF)}\\
\adjustbox{max width=\textwidth}{\input{tables/eq22_usd_gcf.tex}}

\vspace{0.15cm}
\scriptsize\textbf{Out-of-sample (GCF$^{oos}$)}\\
\adjustbox{max width=\textwidth}{\input{tables/eq22_usd_gcf_oos.tex}}
\end{frame}

\begin{frame}{Eq.\ (23) -- USD investor on FXGCF}
\centering
\scriptsize\textbf{In-sample (FXGCF)}\\
\adjustbox{max width=\textwidth}{\input{tables/eq23_usd_fxgcf.tex}}

\vspace{0.15cm}
\scriptsize\textbf{Out-of-sample (FXGCF$^{oos}$)}\\
\adjustbox{max width=\textwidth}{\input{tables/eq23_usd_fxgcf_oos.tex}}
\end{frame}

\begin{frame}{Pooled panel -- Driscoll--Kraay SEs}
\centering
\adjustbox{max width=\textwidth}{\input{tables/panel_dk.tex}}

\vspace{0.2cm}
\small Country fixed effects; DK lag $=\max(18,\,1.3\sqrt{T})$ to handle the 12m return overlap and cross-sectional dependence in GCF.
\end{frame}

\begin{frame}{Out-of-sample: Clark--West}
\centering
\adjustbox{max width=\textwidth}{\input{tables/oos_clark_west.tex}}

\vspace{0.2cm}
\small Expanding window with $\geq 60$ months of training. CW tests the nested-model null that the small ($GCF$-only) model is the DGP; one-sided $p$.
\end{frame}

\begin{frame}{Out-of-sample: GCF vs.\ FXGCF (Diebold--Mariano)}
\centering
\adjustbox{max width=\textwidth}{\input{tables/oos_dm_fxgcf.tex}}

\vspace{0.2cm}
\small Negative DM mean $\Rightarrow$ FXGCF squared OOS errors are smaller than GCF's.
\end{frame}

%---- Fully-OOS factor analysis ---------------------------------
\begin{frame}{Fully out-of-sample factors}
The previous OOS tests used \emph{in-sample} factors as regressors.
Extending the OOS chain so that every factor at $t$ uses only data through $t-1$:
\begin{itemize}
    \item $CF^{oos}_{i,t} = \hat\gamma_{i,1}(1{:}t{-}1)\,c_{i,t}^{(1)}+\hat\gamma_{i,2}(1{:}t{-}1)\,\bar c_{i,t}$
    \item $GCF^{oos}_t = \sum_i w_{i,t}\,CF^{oos}_{i,t}$ \hfill (eq.\ 7 with OOS local CF)
    \item $FXGCF^{oos}_t = \hat\delta_0(1{:}t{-}1) + \hat\delta_1(1{:}t{-}1)\,GCF^{oos}_t$
\end{itemize}

\vspace{0.2cm}
This removes the look-ahead bias that the in-sample $\hat\gamma_{i,j}$ and $\hat\delta_j$ otherwise carry.
\end{frame}

\begin{frame}{In-sample vs.\ OOS GCF and FXGCF}
\centering
\includegraphics[width=0.49\textwidth]{gcf_oos_vs_in.pdf}\hfill
\includegraphics[width=0.49\textwidth]{fxgcf_oos_vs_in.pdf}\\[2pt]
\small Burn-in: 120 months for $CF^{oos}$ (and hence $GCF^{oos}$), plus 60 more for $FXGCF^{oos}$. OOS series track the in-sample factors closely thereafter.
\end{frame}

\begin{frame}{OOS factors -- combined view}
\centering
\includegraphics[width=0.95\textwidth]{oos_factors_combined.pdf}
\end{frame}

\begin{frame}{Campbell--Thompson OOS $R^2$ on OOS factors}
\centering
\adjustbox{max width=\textwidth}{\input{tables/oos_R2_factors.tex}}

\vspace{0.2cm}
\small Benchmark: recursive mean of $rx$. Negative values indicate the factor underperforms the historical mean OOS.
\end{frame}

\begin{frame}{OOS $R^2$ across countries}
\centering
\includegraphics[width=0.95\textwidth]{oos_R2_compare.pdf}
\end{frame}

\begin{frame}{Clark--West with OOS factors}
\centering
\adjustbox{max width=\textwidth}{\input{tables/oos_clark_west_factors.tex}}

\vspace{0.2cm}
\small $H_0$: $GCF^{oos}$-only model is the DGP. Reject in favour of $CF^{oos}+GCF^{oos}$ when $t$ exceeds the one-sided critical value.
\end{frame}

\begin{frame}{Clark--West $t$-stats by country}
\centering
\includegraphics[width=0.95\textwidth]{cw_oos_tstat.pdf}\\[2pt]
\small Reference lines at $1.282/1.645/2.326$ (10\%/5\%/1\% one-sided).
\end{frame}

\begin{frame}{Diebold--Mariano: $GCF^{oos}$ vs.\ $FXGCF^{oos}$}
\centering
\adjustbox{max width=\textwidth}{\input{tables/oos_dm_fxgcf_factors.tex}}

\vspace{0.2cm}
\small Negative DM mean $\Rightarrow$ $FXGCF^{oos}$ outperforms $GCF^{oos}$ for the USD investor.
\end{frame}

\begin{frame}{US replication -- CP-2015 vs.\ CP-2005}
\centering
\adjustbox{max width=0.85\textwidth}{\input{tables/us_replication_R2.tex}}

\vspace{0.4cm}
\includegraphics[width=0.95\textwidth]{us_replication_R2.pdf}
\end{frame}

\begin{frame}{$R^2$ across factor specifications}
\centering
\includegraphics[width=0.95\textwidth]{r2_compare.pdf}
\end{frame}

\begin{frame}{Realized vs.\ CF-fitted excess returns}
\centering
\includegraphics[width=0.95\textwidth]{fitted_vs_realized.pdf}
\end{frame}

%==========================================================
\section{Open Points}
%==========================================================

\begin{frame}{Open methodological points}
\begin{enumerate}
    \item \textbf{Small-sample inference.}
    \cite{bauer_robust_2018} show that standard predictive-regression tests over-reject in small $T$.
    \emph{Plan:} bootstrap inference on $\beta_{CF}$ and $\gamma_{GCF}$; report bias-corrected $p$-values.

    \item \textbf{Look-ahead in the factors.}
    Addressed: $CF^{oos}, GCF^{oos}, FXGCF^{oos}$ are fully recursive.
    \emph{Remaining:} report headline tables on both versions and quantify the look-ahead premium.

    \item \textbf{Multiple testing.}
    Eleven countries $\times$ several specifications inflate FWER.
    \emph{Plan:} headline FDR via Benjamini--Hochberg, plus Bonferroni as a sanity check.

    \item \textbf{Trend inflation specification.}
    $v=0.987,\,M=120$ is calibrated to US survey data.
    \emph{Plan:} robustness with country-specific $v$ and $M\in\{60,180\}$.
\end{enumerate}
\end{frame}

\begin{frame}{Open economic points}
\begin{enumerate}
    \item \textbf{Sample heterogeneity.}
    Pre-/post-GFC and pre-/post-ZLB regimes likely change CF dynamics.
    \emph{Plan:} structural-break diagnostics, sub-sample regressions.

    \item \textbf{Currency-risk decomposition.}
    Why does FXGCF help (or not)?
    \emph{Plan:} decompose $rx^{USD}$ into local $rx$ + FX excess return; attribute the FXGCF gain.

    \item \textbf{Economic value.}
    Statistical $R^2$ vs.\ a portfolio metric \citep{gargano_bond_2019}.
    \emph{Plan:} CER gain from a mean-variance investor timing on CF / GCF / FXGCF.

    \item \textbf{Connection to the literature.}
    Position results against \cite{bauer_interest_2020} and \cite{zhang_global_2021} on long-run trends, and \cite{wright_term_2011} on inflation uncertainty.
\end{enumerate}
\end{frame}

%==========================================================
\section{Plan to Completion}
%==========================================================

\begin{frame}{Plan to completion}
\textbf{05/2026 (now)} -- finish robustness battery (sub-samples, $v$/$M$ grid, OOS-only CF), bootstrap inference, write Section 4 (Results).\\[6pt]
\textbf{06/2026} -- Section 5 (Discussion) tying results to \cite{dahlquist_international_2013}, \cite{bauer_interest_2020}, \cite{wright_term_2011}; economic-value exercise; full draft to supervisor.\\[6pt]
\textbf{07/2026} -- supervisor feedback round, polish, proofread, format, submit.

\vspace{0.4cm}
\textbf{What I would like input on today}
\begin{itemize}
    \item Priority ordering of the open points -- which carry most weight?
    \item Scope of the economic-value section -- is a CER exercise expected?
    \item Presentation choice for the headline result table.
\end{itemize}
\end{frame}

\begin{frame}{Bibliography}
\bibliography{References}
\end{frame}

%==========================================================
\section{Appendix}
%==========================================================

\begin{frame}{G10 sample}
\begin{itemize}
    \item Belgium
    \item Canada
    \item France
    \item Germany
    \item Italy
    \item Japan
    \item Netherlands
    \item Sweden
    \item Switzerland
    \item United Kingdom
    \item United States
\end{itemize}
\end{frame}

\begin{frame}{Trend inflation -- DMA filter (eq.\ 3)}
\[
\pi^{e}_{i,t} = \frac{1-v}{1-v^{M}}\sum_{j=0}^{M-1} v^{j}\,\pi_{i,t-j}
\]
\centering
\includegraphics[width=0.95\textwidth]{trend_panel.pdf}
\end{frame}

\begin{frame}{GDP weights $w_{i,t}$ (eq.\ 8)}
\centering
\includegraphics[width=0.95\textwidth]{gdp_weights.pdf}
\end{frame}

\begin{frame}{FXGCF vs.\ GCF}
\centering
\includegraphics[width=0.95\textwidth]{fxgcf_vs_gcf.pdf}
\end{frame}

\begin{frame}{Pairwise CF correlations across G10}
\centering
\includegraphics[width=0.85\textwidth]{cf_corr_heatmap.pdf}
\end{frame}

\begin{frame}{$t$-statistics by country and spec}
\centering
\includegraphics[width=0.95\textwidth]{tstat_panel.pdf}
\end{frame}

\begin{frame}{Residual ACF: $rx\sim CF$}
\centering
\includegraphics[width=0.95\textwidth]{residual_acf_panel.pdf}\\
\small 12-month overlap induces ACF up to $\sim$ lag 12.
\end{frame}

\begin{frame}{In-sample vs.\ OOS local CF}
\centering
\includegraphics[width=0.95\textwidth]{cf_oos_vs_in.pdf}
\end{frame}

\end{document}
